\section{Statistical Methods}
\label{sec:bg-stats}

The analyses in Chapter~\ref{chap:results} use a small set of standard
procedures, collected here so that the results chapter can apply them without
digression. Observations are \emph{paired}: every arm is evaluated on the
same pull requests, and scores are bounded counts rather than draws from a
normal distribution.

\paragraph{Agreement.}
Cohen's $\kappa$ compares observed agreement between raters with the agreement
expected by chance, giving $1$ for perfect and $0$ for chance-level
agreement~\cite{cohen1960kappa}. Values are read against the bands of Landis
and Koch ($0.41$--$0.60$ moderate, $0.61$--$0.80$ substantial, above $0.80$
almost perfect)~\cite{landiskoch1977}.

\paragraph{Confidence intervals.}
Because the score distributions are neither normal nor of a known form,
intervals are obtained by bootstrapping: the sample is resampled with
replacement, the statistic recomputed on each resample, and the interval read
off the percentiles~\cite{efron1993bootstrap}.

\paragraph{Testing a paired difference.}
A paired permutation test flips the sign of each difference at random and
recomputes the mean; the $p$-value is the proportion of permutations producing
a mean at least as extreme as observed. Where an ordinal reading is
appropriate, the Wilcoxon signed-rank test is used
instead~\cite{wilcoxon1945}.

\paragraph{Effect size.}
For paired data, $d_z$ is the mean difference divided by the standard
deviation of the differences (small at $0.2$, medium at $0.5$, large at
$0.8$). Both $p$-values and effect sizes are reported throughout, since with
thirty-five observations a real but small effect and a null result can be hard to
tell apart from the $p$-value alone.

\paragraph{Multiple comparisons.}
Holm's sequentially rejective procedure controls the family-wise error rate
with more power than Bonferroni by ordering $p$-values and comparing each
against a relaxing threshold~\cite{holm1979}. Chapter~\ref{chap:results}
states for each family which correction was applied and over how many
hypotheses.
